\documentclass[
	%parspace, % Add vertical space between paragraphs
	%noindent, % No indentation of first lines in each paragraph
	%nohyp, % No hyphenation of words
	%twoside, % Double sided format
	%draft, % Quicker draft compilation without rendering images
	%final, % Set final to hide todos
]{elteikthesis}[2026/04/26]

% The minted package is also supported for source highlighting
% See elteikthesis_minted.tex for example
%\usepackage[newfloat]{minted}

% Document's metadata
\title{Source Code Translation Between Programming Languages Using Deep Learning} % title
\date{2026} % year of defense

% Author's metadata
\author{Tamás Fodor}
\degree{Computer Science MSc}

% Superivsor(s)' metadata
\supervisor{Balázs Szalontai} % internal supervisor's name
\affiliation{Assistant Lecturer} % internal supervisor's affiliation
%\extsupervisor{Jane Doe} % external supervisor's name
%\extaffiliation{Senior Developer} % external supervisor's affiliation

% University's metadata
\university{Eötvös Loránd University} % university's name
\faculty{Faculty of Informatics} % faculty's name
\department{Dept. of Software Technology and Methodology} % department's name
\city{Budapest} % city
\logo{elte_cimer_szines} % logo

% Add bibliography file
\addbibresource{elteikthesis.bib}

% The document
\begin{document}

% Set document language
%\documentlang{hungarian}
\documentlang{english}

% List of todos (not in the final document)
%\listoftodos[\todolabel]

% Title page (mandatory)
\maketitle
% Topic declaration page (mandatory) - can also be attached instead
%\includepdf{topicdeclaration.pdf}

% Table of contents (mandatory)
\tableofcontents
\cleardoublepage

\chapter{Introduction}
\label{ch:intro}

%SzB: az intro ne rögtön subsectionnel kezdődjön. Legyen bevezető szöveg. Ez
%minden chapterre igaz. Meg a section-ökre is: ha van section, ne kezdődjüön
%subsectionnal rögtön, stb...

% Main content
\chapter{Introduction}
\label{ch:intro}

Modern software systems are becoming increasingly complex, often consisting of
multiple components written in different programming languages. This diversity
is driven by the need to leverage the strengths of various technologies,
frameworks, and ecosystems. As a consequence, developers frequently face the
challenge of translating or migrating code between programming languages to
ensure maintainability, scalability, and long-term sustainability of software
systems.

In recent years, this problem has increasingly been addressed using deep
learning, particularly Large Language Models (LLMs). These models are capable
of understanding programming languages as structured text and can automatically
generate or translate code between languages based on learned patterns
\cite{llm_limitations}. Unlike traditional rule-based or compiler-based
approaches, LLM-based methods rely on large-scale training data and can capture
both syntactic and semantic aspects of code, making them a promising solution
for automatic code translation \cite{aljagthami2025prompt}.

\section{Background and Motivation}

One of the main motivations behind code translation is the modernization of
legacy systems. Many organizations still rely on outdated programming languages,
which are difficult to maintain and integrate with modern technologies.
Translating such systems reduces technical debt and enables reuse across
platforms \cite{yang2024llm, bruneel2025pipeline}.

Traditional approaches rely on rule-based and compiler-driven techniques, including grammar-based parsing
and abstract syntax tree (AST) transformations \cite{antlr, transpiler_survey,
traditional_translation_survey}. While precise, these methods are difficult to
maintain and scale to complex systems.

Recent advances in deep learning have introduced data-driven approaches that
learn translation patterns from large-scale datasets
\cite{code_translation_survey, llm_survey}. The introduction of the Transformer
architecture \cite{vaswani2017} and LLMs has significantly improved performance in code-related tasks
\cite{brown2020gpt3}. Code-specific models such as CodeLlama \cite{codellama} and
DeepSeek-Coder \cite{deepseek_coder} further enhance performance in programming tasks.

Despite these improvements, automatic code translation remains challenging.
Programming languages require strict syntactic correctness and precise semantic
preservation. Even small mistakes may lead to compilation failures, runtime
errors, or incorrect program behavior. Existing studies show that large language
models often generate code that appears plausible but still contains logical or
semantic errors \cite{llm_limitations, correctness_llm}. Models also struggle
when code depends on larger project context, external libraries, or multiple
source files \cite{evaluation_limitations, long_context_llm}.


Even existing traditional Java-to-C++ translators such as J2C \cite{j2c_tool} frequently require manual post-processing. J2C
attempts to convert Java code into compilable C++ code, but the generated output
still depends on manual handling of Java runtime libraries and missing
dependencies \cite{traditional_translation_survey}. More generally,
transcompilers often preserve the structure of the original code, but the
resulting translations may not follow target-language conventions and typically
require additional manual corrections before they can compile or execute
correctly \cite{roziere2020transcoder}.

Another important challenge is the evaluation of AI-based code translation models.
Traditional metrics such as BLEU \cite{bleu_metric} measure token-level similarity between
generated and reference code, but do not indicate whether the generated program
is actually correct or executable \cite{tran2019ruby}.
Multiple studies have shown that models with high BLEU scores may still fail to
compile or produce incorrect outputs during execution \cite{evtikhiev2022bleu, huang2022exeds}. Because of this, recent research increasingly relies on
execution-based evaluation methods, where generated programs are compiled and
tested against predefined test cases \cite{codescore2023}. These
methods provide a more realistic measure of functional correctness, although
they are also more computationally expensive and dependent on the quality of the
available test cases.




\section{Problem Statement}

%szb: ez nekem túl ismétlős, lehetne kicsit lényegretörőbb a section bevezetése
Recent advances in large language models have significantly improved the
performance of automatic code translation systems, especially for common
programming patterns and widely used language pairs. However, while these models
often produce syntactically correct and readable code, their functional
correctness is not always guaranteed, particularly in more complex scenarios
\cite{llm_limitations, correctness_llm}.

%szb: kicsit fura a sorrend. Az előbb a kiértékelés hangzott másodlagos topic-ként, míg most ez a central
%sztem ez így nemigen jó + az evaluation-t nem is igazán említeném itt, de biztos nem ennyire... Hiszen nem új benchmarkot mutatsz be.
A central challenge in this domain is the evaluation of generated code. Widely
used metrics such as BLEU measure token-level similarity between the generated
and reference code \cite{bleu_metric}, but they do not reflect whether the code
actually compiles or behaves correctly. CodeBLEU improves upon this by
incorporating syntactic and semantic features, but it still does not fully
capture functional correctness \cite{ren2020codebleu}.

To address this issue, execution-based evaluation methods have been introduced,
where generated programs are compiled and tested using predefined test cases.
These methods provide a more realistic measure of correctness, but are more
complex and computationally expensive to apply \cite{huang2022exeds}.

In addition to evaluation challenges, the impact of fine-tuning on code
translation performance is not fully understood. While fine-tuning is commonly
used to adapt pre-trained models to specific tasks, its effectiveness depends on
multiple factors, including the dataset, model architecture, and evaluation
method. %szb: nem igaz, az effectiveness nem múlik a kiértékelési method-on. Helyette az igaz, hogy: az evaluation methdo becsli az effectiveness-t
Therefore, systematic comparison between base and fine-tuned models is
necessary. %szb: esetleg mások csináltak ilyesmit? Miben több/jobb/más a te munkád, mint másoké esetleg? Lehet említeni konkrétumokat
% pl. sokan nem használnak teszt-alapú kiértékelőt, pl. sokan nem finetuneolnak, pl. sokan más adathalmazon finetuneolnak, ilyenek.

This thesis focuses on the comparative evaluation of code translation models. It
investigates how fine-tuning affects model performance, and compares different
models using both similarity-based metrics (BLEU) and execution-based evaluation
methods such as CodeEvalBench. %szb: még nem jutottam a végére, de: such as nem lehet itt. mindent le kell írni, nem csak példákat. ha meg ez az egyedüli ilyen evaluation, akkor a such as használata helytelen
The goal %szb: lehetne ink. úgy fogalmazni, hogy "az eddigieken felül az is érdekes (uninvestigated) kérdés, hogy ezek a mérések hogy viszonyulnak egymáshoz"
is to determine whether improvements observed in traditional metrics correspond
to actual functional correctness.


\section{Objectives and Research Questions}

The main objective of this thesis is to investigate the effectiveness of deep
learning-based approaches for automatic source code translation, with a
particular focus on the effect of fine-tuning and the reliability of different
evaluation methods. Although modern large language models already perform well
on many code-related tasks, it is still unclear to what extent fine-tuning can
improve translation quality in terms of both syntactic accuracy and functional
correctness. %szb: ez miért unclear? indokolni kéne, pl. kevés study van róluk, kevés benchmark, vagy ilyesmik - és persze beciteolni azokat, amik vannak, vagy mégjobb, beciteolni azt, hogy más unclear-nak nevezi ezt

To address this problem, the thesis compares base and fine-tuned versions of %szb: lehetne egyértelműség kedvéért kicsit explicitebb, hogy miben hasonlítod őket össze
multiple code-oriented language models under the same experimental conditions.
The models are fine-tuned on the Java-to-C++ portion of the XLCoST dataset %szb: cite lemaradt 
and evaluated on the CodeEditorBench benchmark %szb: cite lemaradt
using consistent preprocessing, prompting, and evaluation procedures. This setup
makes it possible to compare different model architectures as well as the impact %szb: itt nem kéne már picit explicitebbnek lenni? nem tudom...
of fine-tuning in a fair and reproducible way.

A further objective is to investigate the relationship between similarity-based
metrics and execution-based evaluation methods. Previous studies have shown that
metrics such as BLEU often correlate poorly with functional correctness,
motivating the use of execution-based evaluation \cite{evaluation_limitations,
huang2022exeds}. This thesis therefore examines both types of evaluation in
order to obtain a more complete picture of code translation quality. %szb: le kéne írni, hogyan történik az examination

Based on these objectives, the thesis is guided by the following research
questions:

%szb: amúgy most elgondolkoztam... az RQ1 meg 2 az nem ugyanaz?
% hisz egy translation akkor helyes, ha megőrzi a működést. A modell translation effectiveness-t pedig úgy mérjük, hogy ez hánszor történik
\begin{itemize}
    \item \textbf{RQ1: How effective are deep learning models in translating %szb: lehetne konkretizálni, hogy llm, hisz az van fókuszban
    source code between programming languages?}

    \item \textbf{RQ2: To what extent do these models preserve the functional %szb: nem tudom, hogy a preserve biztos jó szó-e... Esetleg transfer behaviour? vagy ilyesmi? am maradhat ha nincs jó ötlet
    correctness of the original code?} 

    \item \textbf{RQ3: How effective are execution-based evaluation methods
    compared to traditional similarity-based metrics?} %szb: nem jó sztem a szóhasználat: effective a modellekre használandó sztem, itt inkább az a kérdés, hogy:
    %a tradicionális metrikák mennyire jók, illetve hogyan viszonyulnak az execution-based metrikákhoz?
    %kb ez, nem? nem biztos

    \item \textbf{RQ4: What types of errors are most commonly produced by deep
    learning models during code translation?}
\end{itemize}
%ja, és hiányolom a finetuning aspektusát is a rq-kból
%kicsit szkeptikus vagyok ezekkel a kérdésekkel kapcsolatban, bocsána,t hogy csak most írom. javaslok alternatívákat:
% 1 - how effective are current small language models in code translation
% 2 - to what extent can fine tuning improve their translation abilities
% 3 - What types of errors are most commonly produced
% 4 - How well do traditional metrics (such as bleu) estimate the effectiveness of translation llms and how do they relate to execution-based evaluation


\section{Research Questions}

%SzB: Nagyon ismételgeted magad. Az intronak egy tömör pároldalas szövegnek kéne lennie

Building on the challenges discussed in the previous sections, this thesis aims
to investigate the effectiveness and limitations of deep learning-based
approaches for automatic code translation. The focus is placed on evaluating
model performance in realistic scenarios, with particular attention to
functional correctness and reliable evaluation methodologies.

Although recent advances in deep learning have shown promising results in code
translation, existing approaches still face significant challenges. In
particular, models often struggle to generate reliable outputs for complex or
real-world codebases \cite{llm_limitations}. Additionally, traditional
evaluation metrics based on textual similarity are insufficient to assess the
true correctness of translated programs, as they fail to capture behavioral
equivalence \cite{evaluation_limitations}.

To address these challenges, the thesis is guided by the following research
questions:

\begin{itemize}
    \item \textbf{RQ1: How effective are deep learning models in translating
    source code between programming languages?} \\
    This question evaluates the overall performance of modern models, including
    their ability to generate syntactically correct and executable code.

    \item \textbf{RQ2: To what extent do these models preserve the functional
    correctness of the original code?} \\
    This question focuses on whether the translated code maintains the same
    behavior as the original program when evaluated using test cases.

    \item \textbf{RQ3: How effective are execution-based evaluation methods
    compared to traditional similarity-based metrics?} \\
    This question examines whether execution-based approaches provide a more
    reliable assessment of correctness than metrics such as BLEU or token-level
    similarity.

    \item \textbf{RQ4: What types of errors are most commonly produced by deep
    learning models during code translation?} \\
    This question analyzes typical error categories, including syntactic errors,
    semantic inconsistencies, and missing dependencies.
\end{itemize}

These research questions establish the foundation for the experimental and analytical components of the thesis, guiding the evaluation of model performance and the identification of future research directions.

\textcolor{red}{
\section{Structure of the Thesis}


This thesis is structured to provide a clear progression from theoretical foundations to practical implementation and evaluation.

The content of each chapter is organized as follows:

\begin{itemize}

    \item \textbf{Chapter 1} introduces the topic, including the background, motivation, problem statement, objectives, and research questions.

    \item \textbf{Chapter 2} presents the theoretical background, reviewing existing code translation approaches with a focus on deep learning and transformer-based models.

    \item \textbf{Chapter 3} describes the methodology, including the selected datasets, model architecture, training process, and evaluation setup.

    \item \textbf{Chapter 4} focuses on the implementation, covering the development environment, data processing pipeline, model fine-tuning, and integration with evaluation tools.

    \item \textbf{Chapter 5} presents the experimental results and their analysis, including performance evaluation, comparison with baseline methods, and identification of common error patterns.

    \item \textbf{Chapter 6} concludes the thesis by summarizing the key findings, discussing limitations, and outlining directions for future research.
\end{itemize}



\cleardoublepage
}


\chapter{Related work}
\label{ch:rel_work}


\section{Traditional Approaches to Code Translation}

Before the rise of data-driven methods, code translation was primarily addressed using rule-based and compiler-driven techniques. Rule-based approaches rely on manually defined transformation rules that map constructs between programming languages. These systems are typically based on pattern matching and linguistic or syntactic rules, ensuring deterministic and interpretable outputs. However, creating and maintaining such rules requires significant manual effort and does not scale well to complex or evolving codebases \cite{rule_based_translation}.

Another class of traditional solutions includes compiler-based approaches and transpilers. These methods convert source code from one language to another using intermediate representations such as abstract syntax trees (ASTs) or compiler-level structures. While transpilers can provide reliable translations for well-defined patterns, they often depend on handcrafted rules and syntactic mappings, which limits their flexibility and may result in unnatural or suboptimal code \cite{transpiler_limitations}. Furthermore, handling semantic differences between programming languages remains a significant challenge.

In addition, static analysis-based approaches have been used to support code translation by examining program structure, control flow, and dependencies. These techniques aim to preserve correctness by analyzing the behavior of the source program before generating the target code. Despite their advantages, static methods struggle to capture all semantic nuances of modern software systems and are often insufficient when dealing with complex or highly dynamic code \cite{static_analysis_translation}.

In general, traditional approaches provide precise and predictable results but lack scalability and generalizability capabilities. These limitations have led to the increasing adoption of learning-based methods, which are discussed in the following sections.


\section{Deep Learning for Code Translation}

%SzB: Ez inkább theoretical foundations jelleg. Ami nem rossz egyáltalán, sőt
% Lehet át lehetne nevezni a címet theoretical foundations and related work-re
% De a related work rész innen hiányzik. Ha már kódfordítás, akkor itt fel kéne tüntetni, hogy pl. akkor RNN-eket ki használt fordításra. Meg s többbit mindent

The emergence of deep learning has significantly transformed the field of code
translation by enabling data-driven approaches that learn patterns directly from
large-scale datasets. These methods are largely inspired by advances in neural
machine translation (NMT), where sequence-to-sequence (seq2seq) models have
become a dominant paradigm.

Early deep learning approaches to code translation relied on encoder–decoder
architectures based on recurrent neural networks (RNNs), including Long
Short-Term Memory (LSTM) and Gated Recurrent Unit (GRU) models. In these
architectures, the encoder processes the input sequence and encodes it into a
fixed-length vector representation, which is then used by the decoder to
generate the target sequence \cite{seq2seq_original}. While effective in
modeling sequential data, these models suffer from limitations when handling
long sequences, as important information may be lost during encoding
\cite{seq2seq_limitations}. 

To address these limitations, the attention mechanism was introduced, allowing
the model to dynamically focus on relevant parts of the input sequence during
decoding. This significantly improved translation quality by overcoming the
bottleneck of fixed-length representations and enabling better handling of
long-range dependencies \cite{attention_bahdanau}. 

These neural machine translation techniques were later adapted to source code,
where programs are treated as sequences of tokens similar to natural language.
This adaptation enabled models to learn mappings between programming languages
without relying on handcrafted rules. However, unlike natural language, source
code has strict syntax and semantics, which introduces additional challenges for
deep learning models.

Despite the improvements brought by attention mechanisms, RNN-based seq2seq
models remain difficult to parallelize and scale efficiently. This limitation
led to the development of the Transformer architecture, which relies entirely on
attention mechanisms and eliminates recurrence, enabling more efficient training
and improved performance in sequence modeling tasks \cite{vaswani2017}. 

The introduction of Transformers marked a major shift in the field and laid the
foundation for modern approaches to code translation, which are discussed in the
following section.

\section{Transformer-Based Models for Code Understanding and Generation}

%Itt is hiányzik a related work, meg pár hivatkozás, pl. a qwen coderekhez

The introduction of the Transformer architecture marked a major breakthrough in
sequence modeling and has become the foundation of modern approaches to code
understanding and generation. Unlike recurrent models, Transformers rely
entirely on self-attention mechanisms, enabling efficient parallel processing
and improved modeling of long-range dependencies \cite{vaswani2017}.

In recent years, several Transformer-based models have been specifically adapted
for programming languages. These models are typically pre-trained on large-scale
code corpora and are capable of capturing both syntactic and semantic aspects of
source code.

One of the earliest influential models is CodeBERT, a bimodal Transformer
trained on both natural language and programming language data. It learns joint
representations of code and text, making it suitable for tasks such as code
search and documentation generation \cite{feng2020codebert}.

GraphCodeBERT extends this approach by incorporating structural information from
code, particularly data flow. By modeling relationships between variables and
their dependencies, it captures deeper semantic information beyond token
sequences \cite{guo2020graphcodebert}.

Another important model is CodeT5, a unified encoder–decoder Transformer
designed for both code understanding and generation. It introduces
identifier-aware pre-training objectives, allowing the model to better capture
the semantics of variable names and improve performance across various tasks
\cite{wang2021codet5}.

More recent advances include large-scale generative models such as StarCoder and
Qwen Coder, which follow a decoder-only architecture similar to large language
models. These models are trained on massive code datasets and demonstrate strong
performance in code generation and completion tasks.

Despite their success, Transformer-based models still face limitations. In
particular, they often struggle with generalization across different datasets
and may fail to fully capture deeper semantic relationships in complex code
structures \cite{saberi2023model}.
\section{Datasets for Code Translation}

The effectiveness of deep learning models for code translation strongly depends on the availability of high-quality datasets. In recent years, several benchmark datasets have been introduced to support research in code understanding and generation, particularly focusing on parallel code corpora.

One of the most widely used benchmarks is CodeXGLUE, a comprehensive dataset collection designed for multiple code intelligence tasks. It includes a diverse set of datasets covering tasks such as code translation, code summarization, and code search, along with standardized evaluation protocols \cite{lu2021codexglue}. 

For code translation tasks, CodeXGLUE provides parallel datasets containing semantically equivalent code snippets across different programming languages. These datasets are typically constructed by aligning code from open-source repositories, enabling supervised learning of translation models. However, the size and diversity of such parallel datasets are still limited, which can affect the generalization ability of models.

To address these limitations, newer datasets such as XLCoST have been proposed. XLCoST (Cross-Lingual Code Snippet Dataset) is a large-scale benchmark that provides fine-grained parallel code across multiple programming languages. It includes aligned code snippets from several languages and supports a wide range of cross-lingual tasks, making it particularly suitable for code translation research \cite{zhu2022xlcost}.

A key advantage of XLCoST is its focus on multilingual and parallel data, which enables models to learn both syntactic transformations and semantic equivalence more effectively. This makes it more representative of real-world code translation scenarios compared to smaller or less diverse datasets.

In addition to these benchmarks, other datasets have been developed for specific tasks or language pairs. Some of these datasets are synthetically generated, providing clean and well-aligned data, while others are derived from real-world repositories. Synthetic datasets offer scalability but may lack realistic complexity, whereas real-world datasets capture practical challenges such as inconsistent coding styles, dependencies, and incomplete context.

In summary, dataset quality and diversity play a crucial role in the development of robust code translation models. While benchmarks such as CodeXGLUE and XLCoST have significantly advanced the field, challenges remain in obtaining large-scale, high-quality parallel datasets that accurately reflect real-world programming environments.

\section{Evaluation Methods for Code Translation}

% Nem értem... Az előbb már írtál a benchmarkokról, de itt is írsz a kiértékelésről... Kicsit nem áll össze a kép, át kéne gondolni

Evaluating the quality of code translation remains a challenging task, as it
requires assessing not only syntactic correctness but also the preservation of
program semantics. Existing evaluation approaches can be broadly divided into
text-based metrics and execution-based methods.

Text-based metrics are derived from natural language processing and measure the
similarity between generated and reference code at the token level. One of the
most widely used metrics is BLEU, which evaluates n-gram overlap between the
generated output and the reference translation \cite{bleu_metric}. Although BLEU
is simple and widely adopted, it is limited in the context of code, as it does
not account for syntactic structure or semantic equivalence.

To address these limitations, CodeBLEU was introduced as an extension of BLEU
specifically designed for code evaluation. CodeBLEU incorporates additional
factors such as abstract syntax tree (AST) structure and data-flow information,
providing a more comprehensive assessment of code quality
\cite{ren2020codebleu}. However, despite these improvements, CodeBLEU still
relies on surface-level similarity and may fail to capture functional
correctness.

Execution-based evaluation methods provide a more reliable alternative by
assessing whether the generated code behaves correctly when executed. These
approaches typically involve compiling the translated code and running it
against a set of test cases. The correctness of the output is then compared to
the expected results, often using metrics such as pass rate or pass@k. This
method directly measures whether the translated code preserves the intended
functionality.

Recent studies have shown that models achieving high scores on text-based
metrics do not necessarily perform well in execution-based evaluations,
highlighting the limitations of similarity-based approaches
\cite{huang2022exeds}. In particular, code with low textual similarity may still
be functionally correct, while code with high similarity may fail to execute
correctly.

Overall, the discrepancy between text-based and execution-based evaluation
highlights a fundamental challenge in code translation: semantic correctness
cannot be reliably captured by token-level similarity alone. This motivates the
increasing use of execution-based evaluation methods, which better reflect
real-world correctness and are therefore more suitable for assessing practical
code translation systems.

\section{Limitations of Current Approaches}

%SzB: Minden egyes limitation-t alá kell támasztani minimum 1-2 hivatkozással. és azok alapján kell ezeket a gondolatoka megfogalmazni.

Despite significant progress in deep learning-based code translation, current
approaches still face several important limitations that affect their
reliability and applicability in real-world scenarios.

One major issue is hallucination and incorrect logic generation. Models can
produce syntactically valid code that appears plausible but fails to correctly
implement the intended functionality. This highlights the gap between
surface-level correctness and true semantic understanding.

Another limitation is the handling of dependencies and external context. Most
models operate on isolated code snippets and do not fully consider external
libraries, project structure, or runtime dependencies. As a result, generated
code may fail to compile or execute correctly in practical environments.

Context handling also remains a challenge. Although Transformer-based models can
capture long-range dependencies, they are still constrained by fixed context
window sizes. This limitation can lead to incomplete translations or loss of
critical information in larger or more complex codebases.

Furthermore, generalization across datasets and programming languages remains
limited. Models trained on specific datasets often perform well within the same
distribution but struggle when applied to unseen languages, coding styles, or
problem domains.

These limitations are closely connected to challenges in evaluation. Studies
have shown that models achieving high scores on text-based metrics may still
fail in execution-based evaluation, indicating that token-level similarity does
not reliably reflect functional correctness \cite{huang2022exeds}.

These issues demonstrate that current approaches are not yet sufficiently robust
for fully automated code translation. Addressing these challenges requires
improvements in both modeling techniques and evaluation methodologies.

\section{Summary and Research Gap}

The previous sections have reviewed the main approaches to code translation,
including traditional methods, deep learning-based techniques, and Transformer-
based models. These approaches have significantly improved the ability of models
to learn syntactic and semantic patterns from large-scale datasets.

However, several important challenges remain. One key limitation is that current
models often fail to preserve functional correctness, particularly in complex or real-
world scenarios. Additionally, the availability of high-quality parallel datasets is still
limited, which restricts the generalization capabilities of models.

Another major challenge lies in evaluation. Widely used metrics such as BLEU
and CodeBLEU focus primarily on surface-level similarity between generated
and reference code. While CodeBLEU incorporates syntactic and semantic features, it still does not directly evaluate functional correctness \cite{ren2020codebleu}.

Furthermore, studies have shown that these similarity-based metrics do not reli-
ably correlate with human judgment or actual correctness, highlighting the limita-
tions of current evaluation approaches \cite{evtikhiev2022bleu}.

These limitations reveal a clear research gap: there is a lack of reliable evaluation
methodologies that can accurately assess whether translated code is functionally
correct.

This thesis addresses this gap by focusing on execution-based evaluation. By
compiling and testing generated code against test cases, the proposed approach
aims to provide a more reliable assessment of semantic correctness compared to
traditional similarity-based metrics.

%SzB: Nem szoktak külön összefoglalót a related workhöz írni, meg nem így szokták összefoglalni a limitationöket
% Amit itt írok, a 2_6 és 2_7re egyaránt vonatkozik
% A chapter nagyjából ok, de ez a két utolsó rész be kell olvadjon a többibe
% Ahol bemutatsz valamit aminek vannak limitációi, ott lehet megemlíteni félmondatban
% Az is szokott lenni, hogy:
% X et al Y-t csinált, ami ilyen és olyan volt, ugyanakkor Z szempontból nem volt tökéletes
% A et al ezen javított azzal, hogy B-t használt. De persze ez se teljesen ok, mert C.
% És akkor ha itt hirtelen abbahagyod, azzal implicit sejteted, hogy a te munkád amúgy erre fog megoldást kínálni.  

\cleardoublepage

% Bibliography (mandatory)
\phantomsection
\addcontentsline{toc}{chapter}{\biblabel}
\printbibliography[title=\biblabel]
\cleardoublepage

% List of figures (optional) - useful over 3-5 figures
\phantomsection
\addcontentsline{toc}{chapter}{\lstfigurelabel}
\listoffigures
\cleardoublepage

% List of tables (optional) - useful over 3-5 tables
\phantomsection
\addcontentsline{toc}{chapter}{\lsttablelabel}
\listoftables
\cleardoublepage

% List of algorithms (optional) - useful over 3-5 algorithms
\phantomsection
\addcontentsline{toc}{chapter}{\lstalgorithmlabel}
\listofalgorithms
\cleardoublepage

% List of codes (optional) - useful over 3-5 code samples
\phantomsection
\addcontentsline{toc}{chapter}{\lstcodelabel}
\lstlistoflistings
\cleardoublepage

% List of symbols (optional)
%\printnomenclature




\end{document}
